Para el c\'alculo de integrales en regiones que son la imagen de un rect\'angulo bajo una transformaci\'on lineal, como paralelogramos se utiliza un cambio de coordenadas lineal. Dados dos vectores $\mathbf{a}=(a_1,a_2), \mathbf{b}=(b_1,b_2) \in \Rn{2}$ linealmente independientes, un punto $(x,y)$ se relaciona con $(u,v)$ a trav\'ez de la transformaci\'on:

\begin{equation} \label{def:coord_lineales}
\mathbf{T}:D^*\subseteq\Rn{2} \to D\: : \: \mathbf{T}(u, v ) = u\,\mathbf{a} + v\,\mathbf{b} = (a_1 u + b_1 v,\, a_2 u + b_2 v).
\end{equation}

\vspace{0.2cm}

Para el c\'alculo del \emph{Determinante Jacobiano} de $T$ primero calculamos su matriz diferencial $D_\mathbf{T}$, siendo, en este caso, $T$ una transformaci\'on lineal, su matriz diferencial coincide con la matriz de la transformaci\'on,
       \begin{equation*}
       \boldsymbol{D}_{\mathbf{T}}(u, v)
       =
       \begin{pmatrix}
       a_1 & b_1 \\
       a_2 & b_2
       \end{pmatrix}
      \end{equation*}
   luego, el jacobiano es
       \begin{equation}\label{ejemplo_jacobiano_lineales}
       J_{\mathbf{T}}(u, v) = \det(\boldsymbol{D}_{\mathbf{T}})(u, v) = a_1 b_2 - a_2 b_1,
      \end{equation}

es decir, el jacobiano de una transformaci\'on lineal es constante, no depende del punto $(u,v)$ donde se lo eval\'ue, y coincide con el determinante de la matriz que define a $T$.

\vspace{0.2cm}

\textbf{Visualizaci\'on del sistema de coordenadas lineales:}
\begin{figure}[htbp]
    \centering
    % Llamamos al archivo con el código TikZ extraído
    \tikzexternaldisable
\begin{center}
\begin{tikzpicture}[scale=0.9, >=stealth]

  %%% IZQUIERDA: Dominio D* en espacio (u, v)
  % Relleno del cuadrado D*
  \fill[blue!20] (0, 0) rectangle (1, 1);
  % Borde del cuadrado D*
  \draw[blue!70!black, thick] (0, 0) rectangle (1, 1);

  % Ejes u y v
  \draw[->] (-0.6, 0) -- (2.0, 0) node[right] {$u$};
  \draw[->] (0, -0.6) -- (0, 2.0) node[above] {$v$};

  % Marcas y etiquetas en los ejes
  \draw (1, 0.06) -- (1, -0.06) node[below] {\small $1$};
  \draw (0.06, 1) -- (-0.06, 1) node[left] {\small $1$};
  \node[below left] at (0,0) {\small $0$};

  % Etiqueta D*
  \node[color=blue!80!black, font=\large] at (0.5, 0.5) {$D^*$};
  \node[below=0.6cm, font=\small] at (0.5, 0) {Espacio $(u,v)$};

  %%% FLECHA DE TRANSFORMACIÓN T
  \draw[->, ultra thick] (2.6, 1) -- node[above, font=\large] {$T$} (4.1, 1);
  \node[below, font=\small] at (3.35, 1) {$T(u,v)=u\mathbf{a}+v\mathbf{b}$};

  %%% DERECHA: Imagen D (Paralelogramo) en espacio (x, y)
  \begin{scope}[shift={(5.4, 0)}]

    % Relleno del paralelogramo D
    \fill[blue!20] (0,0) -- (1.6,0.5) -- (2.1,2.0) -- (0.5,1.5) -- cycle;
    % Borde del paralelogramo D
    \draw[blue!70!black, thick] (0,0) -- (1.6,0.5) -- (2.1,2.0) -- (0.5,1.5) -- cycle;

    % Ejes cartesianos x e y
    \draw[->] (-0.6, 0) -- (2.9, 0) node[right] {$x$};
    \draw[->] (0, -0.6) -- (0, 2.6) node[above] {$y$};

    % Vectores generadores a y b desde el origen
    \draw[-{Stealth}, ultra thick, red!70!black] (0,0) -- (1.6,0.5);
    \draw[-{Stealth}, ultra thick, red!70!black] (0,0) -- (0.5,1.5);
    \node[below right, color=red!70!black] at (1.6,0.5) {$\mathbf{a}$};
    \node[above left, color=red!70!black] at (0.5,1.5) {$\mathbf{b}$};

    % Vértices O y a+b
    \node[below left] at (0,0) {\small $O$};
    \node[above right] at (2.1,2.0) {\small $\mathbf{a}+\mathbf{b}$};

    % Etiqueta D
    \node[color=blue!80!black, font=\large] at (1.3, 0.85) {$D$};
    \node[below=0.6cm, font=\small] at (1.05, 0) {Espacio $(x,y)$};

  \end{scope}

\end{tikzpicture}
\end{center}
\tikzexternalenable
 
    \label{fig:Visualizacion lineales}
    \caption{Puede observarse que el cuadrado $D^*$ en el espacio $(u,v)$ se transforma en un paralelogramo $D$ en el espacio $(x,y)$ bajo la transformaci\'on lineal $T$.}
\end{figure}


\begin{example}
Hallar el \'area del paralelogramo $D$ de v\'ertices $(0,0)$, $\mathbf{a}=(3,1)$, $\mathbf{b}=(1,2)$ y $\mathbf{a}+\mathbf{b}=(4,3)$.
\end{example}

\textbf{Soluci\'on:} El paralelogramo $D$ es la imagen del cuadrado unitario $D^*=[0,1]\times[0,1]$ bajo la transformaci\'on lineal definida en (\ref{def:coord_lineales}) con $\mathbf{a}=(3,1)$ y $\mathbf{b}=(1,2)$, es decir,
\begin{equation*}
\mathbf{T}:D^*=[0,1]\times[0,1] \subseteq\Rn{2} \to D\: : \: \mathbf{T}(u, v ) = u(3,1) + v(1,2) = (3u+v,\, u+2v).
\end{equation*}

Siguiendo la Definici\'on \ref{def:area}, el \'area se calcula como:
\begin{equation*}
A(D) = \iint_D 1\, dA
\end{equation*}

luego, por el Teorema de cambio de variable (\ref{ teo:cambio_variables_r2}) y recordando (\ref{ejemplo_jacobiano_lineales})
$$ \iint_D 1\, dA = \iint_{D^*} |J_\mathbf{T}|\, du\, dv= \int_0^{1} \int_0^1 |3\cdot 2 - 1 \cdot 1|\, du\, dv $$

por \'ultimo, usando el Teorema de Fubbinni (\ref{col:fubini})

\begin{equation*}
\int_0^{1} \int_0^1 5\, du\, dv = \int_0^{1} \Big(  \int_0^1 5\, du\,  \Big) dv  =  \int_0^{1}  5 \, dv = 5.
\end{equation*}
